\documentclass[10pt,a4paper]{amsart}
\usepackage[margin=19mm]{geometry}
\usepackage{amsmath,amssymb,mathtools}
\usepackage[scr=rsfs,cal=euler]{mathalfa}
\usepackage{xfrac}
\usepackage[unicode,hidelinks,bookmarks=false]{hyperref}
\newcommand{\ie}{i.e.\ }
\newcommand{\cf}{cf.\ }
\newcommand{\resp}{respectively\ }
\newcommand{\Subsection}{\subsection}
\providecommand{\nobreakdash}{\nobreak\hbox{-}}
\setlength{\emergencystretch}{2em}
\multlinegap0pt
\usepackage{amssymb}
\usepackage{xcolor}
\usepackage{enumerate}
\usepackage[normalem]{ulem}
\newtheorem{ass}{Assumption}
\newtheorem{lem}{Lemma}
\newtheorem{prop}{Proposition}
\newcommand{\E}{\mathbb{E}}
\newcommand{\PP}{\mathbb{P}}
\newcommand{\WW}{{\mathbb{W}}}
\newcommand{\bE}{\bar \E}
\newcommand{\bOmega}{\bar \Omega}
\newcommand{\bP}{\bar \PP}
\newcommand{\bcF}{\bar \cF}
\newcommand{\cC}{{\mathcal{C}}}
\newcommand{\cE}{{\mathcal{E}}}
\newcommand{\cF}{{\mathcal{F}}}
\newcommand{\cM}{{\mathcal{M}}}
\newcommand{\dd}{{\mathrm{d}}}
\newcommand{\emark}{\vmark^\dagger}
\newcommand{\emarkfield}{\vmarkfield^\dagger}
\newcommand{\emarklaw}{\vmarklaw^\dagger}
\newcommand{\emarkspace}{\vmarkspace^\dagger}
\newcommand{\nmark}{\vmark'}
\newcommand{\pmark}{a}
\newcommand{\qmark}{b}
\newcommand{\rmark}{c}
\newcommand{\rr}{{\mathbb{R}}}
\newcommand{\tU}{\tilde U}
\newcommand{\tX}{\tilde X}
\newcommand{\vip}{\vskip.12cm}
\newcommand{\vmark}{\xi}
\newcommand{\vmarkfield}{\cE}
\newcommand{\vmarklaw}{\pi}
\newcommand{\vmarkspace}{E}
\title{Mean-field theory via dissociated arrays for particle systems interacting through noisy weights}
\author{Nicolas Fournier, Datong Zhou}
\date{Source: arXiv:2606.12135v1 (CC BY 4.0)}
\begin{document}
\maketitle

\noindent\emph{Source and licence.} Mean-field theory via dissociated arrays for particle systems interacting through noisy weights. Version arXiv:2606.12135v1 (CC BY 4.0), distributed by arXiv under the Creative Commons Attribution 4.0 International licence, http://creativecommons.org/licenses/by/4.0/, as recorded in the arXiv licence metadata for this exact version and shown on the abstract page https://arxiv.org/abs/2606.12135v1. Full licence notice: \texttt{licenses/arxiv-2606.12135-CC-BY-4.0.txt}.\par\smallskip

\noindent\textit{Editorial scope.} Counted unit: the article's prop estlambda (source0.tex lines 584--606). The article's complete original proof of it is delivered with it (source0.tex lines 608--736, one \texttt{proof} environment of 129 lines). No mathematics is written, repaired or re-derived: the statement and the proof are the article's own lines, in the article's own notation, and no retained line is substituted or re-flowed.  Retained uncounted as the source-local context the proof needs: lines 140--157; lines 167--174; lines 178--183; lines 205--232; lines 248--266; lines 508--537.  Nothing was omitted from any retained span, so the package carries no omission record.  No retained line contains a citation, so the wrapper carries no bibliography and no bibliography entry.  No retained line contains a figure environment, an image-inclusion command or drawing code, so no image is delivered and the package declares no asset dependency.  Editorial formatting only, in addition: an AMS \texttt{amsart} preamble that restates the article's own declarations and macros used by the retained lines, namely the environments \texttt{ass}, \texttt{lem}, \texttt{prop} on separate counters, together with the standard packages those lines need.  The retained environments print in this document as Ass 1, Lemma 1, Proposition 1 in order of appearance, whereas the article prints the same items with the numbering of its own full document.  The complete native source of the article as distributed in the arXiv e-print for this exact version is bundled unchanged as \texttt{sources/202601/source0.tex}.
\begin{ass}\label{asr}
The following random objects are independent:
\vip
\noindent $\bullet$ an i.i.d. family $(\vmark_i)_{i\geq 1}$ with common law $\vmarklaw$
valued in some measurable space $(\vmarkspace,\vmarkfield)$,
\vip
\noindent $\bullet$ an i.i.d. family $(\emark_{ij})_{i,j\geq 1, i\neq j}$ with common law $\emarklaw$
valued in some measurable space $(\emarkspace,\emarkfield)$,
\vip
\noindent $\bullet$ some i.i.d. $m$-dimensional Brownian motions $(B^i_t)_{t\geq 0, i\geq 1}$,
\vip
\noindent $\bullet$ some i.i.d. $1$-dimensional Brownian motions $(W^{ij}_t)_{t\geq 0, i,j\geq 1, i\neq j}$.
\vip
There are some measurable functions $F:\vmarkspace\to \rr^d$ and
$G:\vmarkspace^2\times{\emarkspace}\to \rr$ such that, for all $i\geq 1$,
$X_0^i=F(\vmark_i)$ and, for all $i,j\geq 1$ with $i\neq j$,
$U_0^{ij}=G(\vmark_i,\vmark_j,\emark_{ij})$.
\end{ass}

\vip
We will assume at least the following moment condition.
\begin{ass}\label{asm}
For $\vmark_1,\vmark_2,\emark_{12}$ independent with $\vmark_1,\vmark_2\sim \vmarklaw$ and $\emark_{12}\sim\emarklaw$, it holds that 
$$
\E[|F(\vmark_1)|^2+(G(\vmark_1,\vmark_2,\emark_{12}))^2]<\infty.
$$
\end{ass}

\begin{ass}\label{asc}
The functions $b:\rr^d\to \rr^d$, $\sigma:\rr^d \to \cM_{d\times m}(\rr)$,
$\gamma:\rr^d\times\rr^d\to \rr^d$,
$\alpha: \rr^d\times\rr^d\times \rr \to \rr$ and $\beta: \rr^d\times\rr^d\times \rr \to \rr$
are bounded and globally Lipschitz continuous, while $\phi:\rr\to\rr$ is globally Lipschitz continuous.
\end{ass}

\begin{lem}\label{mes}
Grant Assumptions~\ref{asr}, \ref{asm} and~\ref{asc} and set $\cC_\ell=C(\rr_+,\rr^\ell)$. There is a measurable map
$$
\Gamma=(\Gamma_t)_{t\geq 0} : \vmarkspace \times \cC_d\times \vmarkspace \times \cC_d\times\emarkspace\times\cC_1 \to \cC_1
$$
such that for any space $(\Omega,\cF,(\cF_t)_{t\ge 0},\PP)$, any triple $(\vmark,\nmark,\emark)$ of independent $\cF_0$-measurable variables with $\vmark,\nmark\sim\vmarklaw$ and $\emark\sim\emarklaw$,
any pair of continuous $(\cF_t)_{t\geq 0}$-adapted
$\rr^d$-valued processes $X=(X_t)_{t\geq 0}$, $Y=(Y_t)_{t\geq 0}$ such that
\begin{equation}\label{condA}
\text{$\forall\; T>0, \; \exists \; A_T>0$ such that $\forall \;0\leq s<t<s+1\leq T+1$,}\quad \E[|X_t-X_s|^2]\leq
A_T (t-s),
\end{equation}
and such that $Y$ satisfies the same condition,
any $1$-dimensional $(\cF_t)_{t\geq 0}$-Brownian motion $W=(W_t)_{t\geq 0}$,
the pathwise unique solution $U=(U_t)_{t\geq 0}$ to
\begin{equation}\label{dfg}
U_t=G(\vmark,\nmark,\emark)+ \int_0^t \alpha(X_s,Y_s,U_s)\dd s + \int_0^t \beta(X_s,Y_s,U_s)\dd W_s
\end{equation}
satisfies $U=\Gamma(\vmark,X,\nmark,Y,\emark,W)$ in the sense that a.s., for all $t\geq 0$,
$U_t=\Gamma_t(\vmark,X,\nmark,Y,\emark,W)$.
One may moreover build $\Gamma$ in such a way that
for all random sextuplet $(\vmark,X,\nmark,Y,\emark,W)$ as above, a.s.,
\begin{equation}\label{prog}
\text{for all $t\geq 0$,} \quad
\Gamma_t(\vmark,X,\nmark,Y,\emark,W)=\Gamma_t(\vmark,(X_{s\land t})_{s\geq 0},\nmark,
Y,\emark,W).
\end{equation}
\end{lem}

\begin{lem}\label{Lambda}
Grant Assumptions~\ref{asr}, \ref{asm} and~\ref{asc} and suppose either that $\phi$ is bounded or that
\begin{equation}\label{m2c}
\sup_{\pmark,\qmark\in \vmarkspace}\int_{\emarkspace}|G(\pmark,\qmark,\rmark)|^2\emarklaw(\dd \rmark)<\infty.
\end{equation}
For any admissible probability measure $\mu$ on $\vmarkspace\times \cC_d$, we use the
pointwise notation, for $t\geq 0$ and $(\pmark,x)\in \vmarkspace\times \cC_d$,
$$
\Lambda_\mu(t,\pmark,x)=\int_{\vmarkspace\times \cC_d} \int_{\emarkspace\times \cC_1}
\phi(\Gamma_t(\pmark,x,\qmark,y,\rmark,w))\gamma(x_t,y_t) \emarklaw(\dd \rmark)\WW(\dd w)
\mu(\dd \qmark,\dd y).
$$

(i) For each admissible pair $(\vmark,X)$, the map
$t\mapsto \Lambda_\mu(t,\vmark,X)$ is a.s. continuous.
\vip
(ii) If $(\vmark,X)$ is admissible with respect to $(\cF_t)_{t\ge 0}$, then
$(\Lambda_\mu(t,\vmark,X))_{t\geq 0}$ is $(\cF_t)_{t\geq 0}$-adapted.
\end{lem}

\begin{prop}\label{estgamma}
Grant Assumptions~\ref{asr}, \ref{asm} and~\ref{asc}. For all $T>0$, there is
$C_T>0$, depending only on $T,\alpha,\beta$, such that the
following estimates hold.
\vip
(i) For every pair of independent admissible inputs $(\vmark,X)$ and $(\nmark,Y)$ defined on a probability space $(\Omega,\cF,\PP)$,
every auxiliary edge-mark probability space
$(\Omega^\dagger,\cF^\dagger,\PP^\dagger)$ carrying an
$\emarkspace$-valued random variable $\emark$ with law
$\emarklaw$, and every auxiliary Brownian space
$(\Omega^\#,\cF^\#,(\cF_t^\#)_{t\geq0},\PP^\#)$ carrying a one-dimensional
Brownian motion $W$, all independent of each other,
\begin{align*}
\E^\dagger\Big[\E^\#
\Big[\sup_{t\in[0,T]}|\Gamma_t(\vmark,X,\nmark,Y,\emark,W)|^2\Big]
\Big]\leq 2\E^\dagger\big[|G(\vmark,\nmark,\emark)|^2\big]+C_T.
\end{align*}
\vip
(ii) For every pair of vertex-side triples $(\vmark,X,\tilde X)$ and
$(\nmark,Y,\tilde Y)$ such that $(\vmark,X)$, $(\vmark,\tilde X)$, $(\nmark,Y)$ and
$(\nmark,\tilde Y)$ are admissible and such that $(\vmark,X,\tilde X)$ is independent of $(\nmark,Y,\tilde Y)$, and every pair of auxiliary spaces as in (i),
independent of these triples, for all $t\in[0,T]$,
\begin{align*}
&\E^\dagger\Big[\E^\#
\Big[\sup_{u\in[0,t]}|\Gamma_u(\vmark,X,\nmark,Y,\emark,W)
-\Gamma_u(\vmark,\tilde X,\nmark,\tilde Y,\emark,W)|^2\Big]\Big]\\
\leq& C_T\Big(\sup_{s\in[0,t]}|X_s-\tilde X_s|^2
+\sup_{s\in[0,t]}|Y_s-\tilde Y_s|^2\Big).
\end{align*}
\end{prop}

\begin{prop}\label{estlambda}
Grant Assumptions~\ref{asr}, \ref{asm} and~\ref{asc} and suppose either that
$\phi$ is bounded or that~\eqref{m2c} holds.
For all $T>0$, there is $C_T>0$ such that the following estimates hold.
\vip
(i) For any admissible probability measure $\mu$ on
$\vmarkspace\times\cC_d$, any admissible pair $(\vmark,X)$, any
$t\in[0,T]$,
\begin{equation*}
|\Lambda_\mu(t,\vmark,X)|\leq C_T\quad \hbox{a.s.}
\end{equation*}
\vip
(ii) Let $(\Omega,\cF,(\cF_t)_{t\geq0},\PP)$ support an $\cF_0$-measurable random
variable $\vmark$ and continuous adapted processes $X=(X_t)_{t\geq0}$ and
$\tX=(\tX_t)_{t\geq0}$ such that $(\vmark,X)$ and $(\vmark,\tX)$ are admissible.
Write $\Lambda_{\vmark,X}:=\Lambda_{\mathrm{Law}(\vmark,X)}$ and
$\Lambda_{\vmark,\tX}:=\Lambda_{\mathrm{Law}(\vmark,\tX)}$. Then, for all
$t\in[0,T]$,
\begin{equation*}
\E\Big[|\Lambda_{\vmark,X}(t,\vmark,X)-\Lambda_{\vmark,\tX}(t,\vmark,\tX)|^2\Big]
\leq C_T\E\Big[\sup_{s\in[0,t]}|X_s-\tX_s|^2\Big].
\end{equation*}
\end{prop}

\begin{proof}
We start with (i). Additionally to $(\Omega,\cF,(\cF_t)_{t\geq0},\PP)$ on which is defined $(\xi,X)$, we introduce a probability space
$(\Omega^*,\cF^*,(\cF_t^*)_{t\geq0},\PP^*)$ endowed with an admissible pair
$(\nmark,Y)$ with law $\mu$, a probability space
$(\Omega^\dagger,\cF^\dagger,\PP^\dagger)$ endowed with an
$\emarkspace$-valued random variable $\emark$ with
law $\emarklaw$, and a probability space
$(\Omega^\#,\cF^\#,(\cF_t^\#)_{t\geq0},\PP^\#)$ endowed with a one-dimensional
Brownian motion $W$. The three auxiliary spaces are taken independent of $(\Omega,\cF,(\cF_t)_{t\geq0},\PP)$ and of each other. Recalling Lemma~\ref{Lambda},
$$
\Lambda_{\mu}(t,\xi,X)=\E^*\Big[\E^\dagger\Big[\E^\#\Big[
\phi(\Gamma_t(\vmark,X,\nmark,Y,\emark,W)) \gamma(X_t,Y_t)
\Big]\Big]\Big].
$$
Define
\begin{equation*}
\Psi_{\mu,T}(\vmark,X):=
\E^*\Big[\E^\dagger\Big[\E^\#\Big[
\sup_{t\in [0,T]} |\phi(\Gamma_t(\vmark,X,\nmark,Y,\emark,W))|
\Big]\Big]\Big].
\end{equation*}
If $\phi$ is bounded, then $\Psi_{\mu,T}(\vmark,X)\leq ||\phi||_\infty$. If~\eqref{m2c} holds, then we use that $|\phi(u)|\leq C(1+|u|)$, Proposition~\ref{estgamma}-(i), together with the Cauchy--Schwarz inequality. This gives \begin{align*}
\Psi_{\mu,T}(\vmark,X)
&\leq C+C\E^*\Big[\E^\dagger\Big[\E^\#
\Big[\sup_{t\in[0,T]}|\Gamma_t(\vmark,X,\nmark,Y,\emark,W)|\Big]\Big]\Big]
\leq C+ C_T \E^*[\E^\dagger[|G(\xi,\xi',\xi^\dagger)|^2]]^{\frac12}.
\end{align*}
This last quantity is bounded by some constant $C_T$ thanks to~\eqref{m2c}.
\vip

Then~(i) follows from
$|\Lambda_\mu(t,\vmark,X)|\leq ||\gamma||_\infty\Psi_{\mu,T}(\vmark,X)$.
\vip
We now prove (ii). Additionally to $(\Omega,\cF,(\cF_t)_{t\geq0},\PP)$, we
introduce a second probability space
$(\Omega^*,\cF^*,(\cF_t^*)_{t\geq0},\PP^*)$ endowed with a copy
$(\vmark^*,X^*,\tX^*)$ of $(\vmark,X,\tX)$, as well as a third probability space
$(\Omega^\dagger,\cF^\dagger,\PP^\dagger)$ endowed with an
$\emarkspace$-valued random variable $\emark$ with
law $\emarklaw$ and a fourth probability space
$(\Omega^\#,\cF^\#,(\cF_t^\#)_{t\geq0},\PP^\#)$ endowed with a one-dimensional
Brownian motion $W$. On
$$
(\bOmega,\bcF,(\bcF_t)_{t\geq0},\bP)
=\big(\Omega\times\Omega^*\times\Omega^\dagger\times\Omega^\#,
\cF\otimes\cF^*\otimes\cF^\dagger\otimes\cF^\#,
(\cF_t\otimes\cF_t^*\otimes\cF^\dagger\otimes\cF_t^\#)_{t\geq0},
\PP\otimes\PP^*\otimes\PP^\dagger\otimes\PP^\#\big),
$$
the four families $(\vmark,X,\tX)$, $(\vmark^*,X^*,\tX^*)$,
$\emark$ and $W$ are independent. We consider the pathwise unique
solutions $U$ and $\tU$ to
\begin{align*}
U_t=&G(\vmark,\vmark^*,\emark)
+\int_0^t \alpha(X_s,X_s^*,U_s)\dd s
+\int_0^t \beta(X_s,X_s^*,U_s)\dd W_s,\\
\tU_t=&G(\vmark,\vmark^*,\emark)
+\int_0^t \alpha(\tX_s,\tX_s^*,\tU_s)\dd s
+\int_0^t \beta(\tX_s,\tX_s^*,\tU_s)\dd W_s.
\end{align*}
By Lemma~\ref{mes},
$U=\Gamma(\vmark,X,\vmark^*,X^*,\emark,W)$ and
$\tU=\Gamma(\vmark,\tX,\vmark^*,\tX^*,\emark,W)$. Thus, as in (i),
$$
\Lambda_{\vmark,X}(t,\vmark,X)
=\E^*\Big[\E^\dagger\Big[\E^\#\big[\phi(\Gamma_t(\vmark,X,\vmark^*,X^*,\emark,W))\gamma(X_t,X_t^*)\big]\Big]\Big]
=\E^*\Big[\E^\dagger\Big[\E^\#\big[\phi(U_t)\gamma(X_t,X_t^*)\big]\Big]\Big].
$$
Similarly,
$\Lambda_{\vmark,\tX}(t,\vmark,\tX)
=\E^*[\E^\dagger[\E^\#[\phi(\tU_t)\gamma(\tX_t,\tX_t^*)]]]$, so that
\begin{align}
\E\Big[|\Lambda_{\vmark,X}(t,\vmark,X)-\Lambda_{\vmark,\tX}(t,\vmark,\tX)|^2\Big]
=& \E\Big[\Big|\E^*\Big[\E^\dagger\Big[\E^\#\Big[\phi(U_t)\gamma(X_t,X_t^*)
-\phi(\tU_t)\gamma(\tX_t,\tX_t^*)\Big]\Big]\Big]\Big|^2\Big]\notag\\
\leq& \E\Big[\E^*\Big[\E^\dagger\Big[\E^\#\Big[
\Big|\phi(U_t)\gamma(X_t,X_t^*)
-\phi(\tU_t)\gamma(\tX_t,\tX_t^*)\Big|^2\Big]\Big]\Big]\Big]\notag\\
=& \bE\Big[\Big|\phi(U_t)\gamma(X_t,X_t^*)
-\phi(\tU_t)\gamma(\tX_t,\tX_t^*)\Big|^2\Big]. \label{lpd}
\end{align}
\vip

{\it Case 1: $\phi$ bounded.} Using that $\phi,\gamma$ are bounded and Lipschitz
continuous, \eqref{lpd} implies that
\begin{align}\label{1jab1}
\E\Big[|\Lambda_{\vmark,X}(t,\vmark,X)-\Lambda_{\vmark,\tX}(t,\vmark,\tX)|^2\Big]
\leq& C\bE[|U_t-\tU_t|^2]
+C\bE[|X_t-\tX_t|^2+|X_t^*-\tX_t^*|^2].
\end{align}
By Proposition~\ref{estgamma}-(ii),
\begin{equation}\label{1jab2}
\bE[|U_t-\tU_t|^2]\leq
C_T\bE\Big[\sup_{s\in[0,t]}\big(|X_s-\tX_s|^2+|X_s^*-\tX_s^*|^2\big)\Big].
\end{equation}
Inserting~\eqref{1jab2} into~\eqref{1jab1}, we find
\begin{align*}
\E\Big[|\Lambda_{\vmark,X}(t,\vmark,X)\!-\!\Lambda_{\vmark,\tX}(t,\vmark,\tX)|^2\Big]
\leq C_T\bE\Big[\sup_{s\in[0,t]}
\big(|X_s\!-\!\tX_s|^2\!+\!|X_s^*\!-\!\tX_s^*|^2\big)\Big]
\leq C_T\E\Big[\sup_{s\in[0,t]}|X_s\!-\!\tX_s|^2\Big],
\end{align*}
as desired, because $(X^*,\tX^*)$ has the same law as $(X,\tX)$.
\vip
{\it Case 2: \eqref{m2c} holds.} Since
$\phi$ has linear growth, since $\phi,\gamma$ are Lipschitz continuous and since $\gamma$ is bounded, \eqref{lpd} implies that
\begin{align}\label{jab1}
\E\Big[|\Lambda_{\vmark,X}(t,\vmark,X)-\Lambda_{\vmark,\tX}(t,\vmark,\tX)|^2\Big]
\leq& C\bE[|U_t-\tU_t|^2]\notag\\
&+C\bE\Big[(1+|U_t|^2)
\big(|X_t-\tX_t|^2+|X_t^*-\tX_t^*|^2\big)\Big].
\end{align}
Moreover, \eqref{1jab2} still holds, while Proposition~\ref{estgamma}-(i) and~\eqref{m2c}
give, uniformly in $(\omega,\omega^*)\in \Omega\times\Omega^*$,
$$
\E^\dagger\big[\E^\#[1+|U_t|^2]\big]\leq C_T.
$$
Consequently,
\begin{align}
&\bE\Big[(1+|U_t|^2)
\big(|X_t-\tX_t|^2+|X_t^*-\tX_t^*|^2\big)\Big]\notag\\
&=\E\Big[\E^*\Big[\big(|X_t-\tX_t|^2+|X_t^*-\tX_t^*|^2\big)
\E^\dagger\big[\E^\#[1+|U_t|^2]\big]\Big]\Big]\notag\\
&\leq C_T\E\Big[\E^*\big[|X_t-\tX_t|^2+|X_t^*-\tX_t^*|^2\big]\Big]\notag\\
&\leq C_T\E\Big[\sup_{s\in[0,t]}|X_s-\tX_s|^2\Big]. \label{jab2}
\end{align}
The conclusion follows by inserting~\eqref{1jab2} and~\eqref{jab2}
into~\eqref{jab1}.
\end{proof}

\end{document}
